% FILE: 06_contribution_to_thesis.tex
% Project: research-prompt-manufactory
% Thesis Role: prompt-manufacturing-and-subagent-command-surface

\section{Contribution to the Dissertation}
\label{sec:research-prompt-manufactory-contribution}

\subsection{Contribution to Prediction vs.~Coordination}
\label{sec:research-prompt-manufactory-contribution-prediction}

The central thesis argument distinguishes prediction (language model inference) from
coordination (multi-agent process execution). The Prompt Manufactory contributes to this
argument by demonstrating that the command surface of a coordination system cannot be
derived from prediction alone.

A hand-written prompt is, in effect, a prediction about what instruction would cause a
subagent to behave correctly. It encodes the author's model of what the subagent should
do, expressed in natural language, and relies on the language model's prediction
machinery to interpret and act on that model. The instruction has no formal relationship
to the process law it purports to enforce. Whether the subagent's actual behavior
conforms to the declared workflow is not derivable from the prompt text.

The Prompt Manufactory manufactures the instruction from the law itself. The subagent
warrant for a Census role (\texttt{role\_engine\_census}) is manufactured from the
\texttt{pm:SubagentRole} instance that carries the formal owned surfaces
(\texttt{src/mining}, \texttt{src/replay}, \texttt{src/conformance}, \texttt{src/ocpq},
\texttt{src/receipt}, \texttt{tests/}), the formal forbidden surfaces
(\texttt{../wasm4pm-compat/}, \texttt{../process-intelligence/}), the output contract
(\texttt{research/pi-program/intel/wasm4pm-engine-census.md}), and the refusal gate
(\texttt{pm:REFUSAL\_MISSING\_SOURCE}). The manufactured warrant is not a prediction
about correct behavior; it is a rendering of the formal definition of correct behavior.

% AUTHOR THESIS / INTERPRETATION
This distinction is the coordination contribution: a manufactured prompt creates a
verifiable alignment between the declared law and the operational instruction, a
property that prediction-derived prompts cannot provide.

\subsection{Contribution to Process Evidence}
\label{sec:research-prompt-manufactory-contribution-evidence}

The Van der Aalst Constitution governing this research foundry requires that process
claims be provable from event evidence. The Prompt Manufactory contributes a
manufacturing receipt chain that constitutes event evidence for the claim that subagent
instructions were lawfully derived.

Each emitted warrant's receipt entry in \texttt{prompt-receipt-ledger.md} records a
manufacturing event with timestamp, source ontology URI, manufacturing method, and
proof gate results. This receipt is queryable evidence that the instruction was
manufactured, not authored. When the ggen v26.5.21 execution manifest records 41
artifacts emitted across 6 rules with receipts-equal-artifacts gate PASS
(\texttt{2026-06-01T23:09:45Z}), this constitutes a manufacturing event log from which
the claim ``all warrants were manufactured from graph law'' can be partially verified.

The contribution is partial because: the fallback manufacture of the proof-specimen
warrant used a manual method (ggen v5 template rendering failed on empty SPARQL result
set), and the receipt correctly records this as \texttt{COMPLETE\_VIA\_FALLBACK} rather
than full automated manufacture. Process evidence that records its own deviations
from the declared process is stronger evidence than evidence that claims conformance
without deviation records.

\subsection{Contribution to Manufacturing Architecture}
\label{sec:research-prompt-manufactory-contribution-architecture}

The Prompt Manufactory establishes that the post-cyberpunk manufacturing layer can
extend to the command surface of the system it manufactures. In the full manufacturing
architecture, the sequence is: law encoded as RDF $\to$ SPARQL extraction $\to$ Tera
rendering $\to$ Markdown artifact $\to$ receipt. This sequence is the same pipeline
used by ggen v26.5.21 to manufacture wasm4pm compatibility laws, OTel registry
artifacts, and ZOEapp projection surfaces. The Prompt Manufactory applies this same
pipeline recursively: the manufacturing system manufactures the instructions for its
own subagents.

The architectural contribution is the \texttt{pm:RenderedPrompt} class with its
required \texttt{pm:derivedFrom} triple and its SHACL enforcement. This creates a
formal warrant type that is distinct from a document, a configuration file, or a
natural-language instruction. A \texttt{pm:RenderedPrompt} is a typed manufacturing
artifact in the same ontological sense as a \texttt{pm:Receipt} or a
\texttt{pm:Checkpoint} --- it carries provenance, type, and proof requirements.

The manufacturing manifest structure (6 rules, 41 artifacts, 5 artifact types, 41,867
bytes, receipts-equal-artifacts gate) demonstrates that warrant manufacturing at
research-program scale is operationally feasible with a single \texttt{ggen sync}
command.

\subsection{Contribution to Capital-Flow Infrastructure}
\label{sec:research-prompt-manufactory-contribution-capital}

% AUTHOR THESIS / INTERPRETATION
The thesis lineage extends to capital-flow, settlement, and DAO infrastructure. The
Prompt Manufactory's contribution to this layer is indirect but structurally important:
it establishes that the subagent command surface of a complex, multi-party system
can be manufactured from formal law rather than authored by individual participants.

In a capital-flow or DAO context, the commands issued to execution agents must be
auditable and derivable from the governing rules of the system. A hand-authored
instruction to an execution agent handling settlement has no formal provenance chain
to the governing rules; a manufactured warrant does. The \texttt{pm:Receipt} class and
its \texttt{pm:proves} link to the \texttt{pm:RenderedPrompt} it certifies establishes
the proof-of-manufacture pattern that capital-flow infrastructure would require to
demonstrate that execution instructions were lawfully derived from governance law.

The Prompt Manufactory does not implement capital-flow infrastructure directly. Its
contribution is to prove the warrant-manufacturing pattern at the research-program
layer, establishing the architecture that would need to be instantiated at the
capital-flow layer for that infrastructure to be auditable.

% FUTURE WORK
The extension of the \texttt{pm:PromptClass} enumeration to include settlement and
governance warrant types (beyond the current INTEL $\mid$ RECONCILE $\mid$ REMEDIATE
$\mid$ SYNC $\mid$ AUDIT $\mid$ CHECKPOINT $\mid$ SPR $\mid$ LIVESTREAM) is deferred
as future work.
